% =====================================================================
% Cylindrical Affinity Kernel — includable section
%
% Usage in parent document:
%   1. Add to your .bib file (or \addbibresource):
%        cylindrical_kernel_refs.bib
%   2. Then: % =====================================================================
% Cylindrical Affinity Kernel — includable section
%
% Usage in parent document:
%   1. Add to your .bib file (or \addbibresource):
%        cylindrical_kernel_refs.bib
%   2. Then: % =====================================================================
% Cylindrical Affinity Kernel — includable section
%
% Usage in parent document:
%   1. Add to your .bib file (or \addbibresource):
%        cylindrical_kernel_refs.bib
%   2. Then: % =====================================================================
% Cylindrical Affinity Kernel — includable section
%
% Usage in parent document:
%   1. Add to your .bib file (or \addbibresource):
%        cylindrical_kernel_refs.bib
%   2. Then: \input{cylindrical_kernel_includable}
%
% Requires: amsmath, amssymb, enumitem (in parent preamble)
% =====================================================================

\section{Adaptive Spatial--Dynamical Affinity Blending for Macro-State Discovery}

\subsection{Problem}

The pendulum state space is parameterised by $(\theta, \omega)$ where
$\theta \in [-\pi, \pi]$ is \emph{periodic}---the states $\theta = -\pi$
and $\theta = +\pi$ are physically identical (pendulum hanging straight
down).  However, in the discretised grid they occupy opposite edges, so
standard spectral clustering on the successor representation matrix
$\mathbf{M}$ can fragment dynamically-connected boundary states into
separate clusters.

\subsection{Cylindrical Embedding}

Instead of operating in the flat $(\theta, \omega)$ space, we embed each
discrete state $s_{i,j}$ (with bin centres $\theta_i$, $\omega_j$) into
a 3-dimensional \emph{cylindrical coordinate space}:
\begin{equation}
\mathbf{x}_{i,j} = \left(\cos\theta_i,\;\sin\theta_i,\;\frac{\omega_j}{\omega_{\max}}\right)
\end{equation}
where $\omega_{\max} = 8.0\;\text{rad/s}$ is the clipping bound for
angular velocity.  The angle maps onto the unit circle
($|\cos\theta|, |\sin\theta| \leq 1$), and dividing $\omega$ by
$\omega_{\max}$ normalises it to $[-1, 1]$, making all three dimensions
commensurate.

In this embedding, states at $\theta = -\pi$ and $\theta = +\pi$ map to
the \emph{same point} $(-1, 0, \omega_j/\omega_{\max})$, eliminating the
boundary discontinuity.

\subsection{Spatial RBF Kernel}

We compute pairwise squared Euclidean distances in the cylindrical space:
\begin{equation}
D^2_{(i,j),(k,l)} = \left\|\mathbf{x}_{i,j} - \mathbf{x}_{k,l}\right\|^2
\end{equation}
and form a radial basis function (Gaussian) kernel:
\begin{equation}
K^{\text{spatial}}_{(i,j),(k,l)} = \exp\!\left(-\frac{D^2_{(i,j),(k,l)}}{2\sigma^2}\right)
\end{equation}
with bandwidth $\sigma = 1.0$.  This kernel is symmetric, non-negative,
and equals~1 on the diagonal.

For non-periodic state spaces (e.g.\ Mountain Car), the embedding reduces
to normalised Cartesian coordinates and the kernel becomes a standard
spatial RBF.

\subsection{SR-Based Affinity}

The standard spectral clustering affinity is the symmetrised successor
matrix:
\begin{equation}
\mathbf{A}^{\text{SR}} = \max\!\left(\mathbf{M},\;\mathbf{M}^\top\right)
\end{equation}
To blend it with the spatial kernel, we first normalise
$\mathbf{A}^{\text{SR}}$ to $[0, 1]$:
\begin{equation}
\widetilde{\mathbf{A}}^{\text{SR}} = \frac{\mathbf{A}^{\text{SR}} - \min\!\left(\mathbf{A}^{\text{SR}}\right)}{\max\!\left(\mathbf{A}^{\text{SR}}\right) - \min\!\left(\mathbf{A}^{\text{SR}}\right)}
\end{equation}

\subsection{Adaptive Blending}

A uniform blend weight $\alpha$ applies the spatial kernel equally to all
state pairs, regardless of how well $\mathbf{M}$ has been estimated.
This can be problematic: in environments with walls or barriers, two
states may be spatially adjacent but dynamically disconnected.  A uniform
spatial kernel would incorrectly boost their affinity, blurring the very
bottleneck structure that clustering should discover.

We address this with an \emph{adaptive} blend weight that varies per
state pair based on the confidence in $\mathbf{M}$.  The intuition is
simple: trust $\mathbf{M}$ where it is reliable; fall back on the spatial
kernel where it is not.

\paragraph{Per-state confidence.}
The row sum $r_i = \sum_j M_{ij}$ reflects how well state~$i$ has been
explored: well-visited states accumulate large successor mass, while
rarely-visited states (e.g.\ near-goal boundary regions) have small row
sums.  We normalise to $[0,1]$:
\begin{equation}
\hat{c}_i = \frac{r_i}{\max_k\, r_k}
\end{equation}

\paragraph{Per-pair confidence.}
For the pair $(i,j)$, we take the minimum of the two states'
confidences, so that the spatial kernel activates whenever \emph{either}
state is poorly estimated:
\begin{equation}
c_{ij} = \min\!\left(\hat{c}_i,\;\hat{c}_j\right)
\end{equation}

\paragraph{Adaptive blend weight.}
The effective spatial weight for pair $(i,j)$ is:
\begin{equation}
\alpha_{ij} = \alpha_{\max}\,(1 - c_{ij})
\end{equation}
where $\alpha_{\max}$ is the maximum blend weight (a hyperparameter;
default $0.5$ for the Pendulum, $0.3$ for Mountain Car).  The final
blended affinity becomes:
\begin{equation}
\boxed{A_{ij} = (1 - \alpha_{ij})\,\widetilde{A}^{\text{SR}}_{ij} + \alpha_{ij}\,K^{\text{spatial}}_{ij}}
\end{equation}

This gives the desired behaviour in three regimes:

\begin{itemize}[leftmargin=2em]
\item \textbf{Both states well-visited}
($\hat{c}_i, \hat{c}_j \approx 1$): $\alpha_{ij} \approx 0$.
The affinity is determined almost entirely by the SR---dynamical barriers
(walls, unreachable regions) are preserved.

\item \textbf{At least one state poorly-visited}
($\min(\hat{c}_i, \hat{c}_j) \approx 0$):
$\alpha_{ij} \approx \alpha_{\max}$.  The spatial kernel fills in for the
unreliable SR, preventing sparse boundary states from splitting off as
outlier clusters.

\item \textbf{Periodic boundaries (Pendulum)}:  States at
$\theta \approx \pm\pi$ are typically well-visited (the pendulum swings
through the bottom frequently), so $c_{ij}$ is moderate-to-high.  The
cylindrical kernel still provides a boost because the \emph{spatial}
distance across the branch cut is zero, while the SR may under-represent
this connection due to bin discretisation.  The adaptive weight ensures
this boost is proportionate.
\end{itemize}

The key property is that $\mathbf{K}^{\text{spatial}}$ is invariant to
the choice of angle branch cut---$\theta = -\pi$ and $\theta = +\pi$
have zero distance---so boundary states that are dynamically connected in
$\mathbf{M}$ are no longer artificially separated by the discretisation.

\subsection{Scope of Application}

The blended affinity $\mathbf{A}$ is used exclusively for \emph{macro-state
discovery}: it determines the spectral embedding (for visualisation) and
the spectral clustering (for cluster assignments).  The core micro-level
quantities---the learned successor matrix $\mathbf{M}$, the transition
matrix $\mathbf{B}$, and the value function
$\mathbf{V} = \mathbf{M}\,\mathbf{C}$---remain unmodified.

Macro-level planning matrices ($\mathbf{B}_{\text{macro}}$,
$\mathbf{M}_{\text{macro}}$) are derived from raw transition counts
between the resulting clusters, so they are \emph{indirectly} shaped by
the blended clustering but do not incorporate the spatial kernel
themselves.

This separation is deliberate: the spatial kernel should produce
physically coherent macro-state regions without distorting the learned
dynamics that drive action selection.

\subsection{Related Work}

Machado et~al.~\autocite{machado2018eigenoption} proved that the
eigenvectors of the successor representation matrix $\mathbf{M}$ are
equivalent to proto-value functions---the eigenvectors of the graph
Laplacian over the state-transition graph.  Their \emph{eigenoption
discovery} framework uses spectral decomposition of $\mathbf{M}$ to
identify bottleneck states and construct temporally abstract actions
(options).  The earlier Laplacian framework~\autocite{machado2017laplacian}
showed that eigenoptions correspond to eigenvectors associated with the
largest eigenvalues of the graph Laplacian, providing a principled
connection between spectral methods and option discovery in reinforcement
learning.  Our micro-level spectral clustering on $\mathbf{M}$ follows
this established approach directly.

Clustering data embedded in spaces with periodic boundary conditions is a
well-studied problem in computational science.  The established solution
is to embed periodic coordinates on the unit circle
$(\cos\theta, \sin\theta)$ so that the metric respects the underlying
topology.  This approach has been applied to
$k$-means~\autocite{gao2022kmeans} and
DBSCAN~\autocite{lazzari2025dbscan} in molecular simulation and other
periodic settings.

The idea of combining multiple affinity matrices appears extensively in
multi-view clustering, where a unified affinity is learned from several
``views'' of the same data.  Constrained spectral clustering modifies the
affinity matrix with must-link / cannot-link constraints or spatial
regularisation terms~\autocite{vonluxburg2007tutorial}.  Our convex
combination of a dynamics-based affinity
($\widetilde{\mathbf{A}}^{\text{SR}}$) with a geometry-based kernel
($\mathbf{K}^{\text{spatial}}$) is conceptually analogous, treating the
successor representation and the physical state layout as two
complementary views of the same state space.

To our knowledge, the specific combination of a learned successor
representation with a cylindrical spatial RBF kernel via an adaptive,
confidence-weighted convex blend for macro-state discovery has not
appeared in prior work.  The adaptive confidence weighting is a new
contribution that makes the blending robust to environments containing
dynamical barriers, where a uniform spatial kernel would incorrectly
smooth over true bottleneck structure.

%
% Requires: amsmath, amssymb, enumitem (in parent preamble)
% =====================================================================

\section{Adaptive Spatial--Dynamical Affinity Blending for Macro-State Discovery}

\subsection{Problem}

The pendulum state space is parameterised by $(\theta, \omega)$ where
$\theta \in [-\pi, \pi]$ is \emph{periodic}---the states $\theta = -\pi$
and $\theta = +\pi$ are physically identical (pendulum hanging straight
down).  However, in the discretised grid they occupy opposite edges, so
standard spectral clustering on the successor representation matrix
$\mathbf{M}$ can fragment dynamically-connected boundary states into
separate clusters.

\subsection{Cylindrical Embedding}

Instead of operating in the flat $(\theta, \omega)$ space, we embed each
discrete state $s_{i,j}$ (with bin centres $\theta_i$, $\omega_j$) into
a 3-dimensional \emph{cylindrical coordinate space}:
\begin{equation}
\mathbf{x}_{i,j} = \left(\cos\theta_i,\;\sin\theta_i,\;\frac{\omega_j}{\omega_{\max}}\right)
\end{equation}
where $\omega_{\max} = 8.0\;\text{rad/s}$ is the clipping bound for
angular velocity.  The angle maps onto the unit circle
($|\cos\theta|, |\sin\theta| \leq 1$), and dividing $\omega$ by
$\omega_{\max}$ normalises it to $[-1, 1]$, making all three dimensions
commensurate.

In this embedding, states at $\theta = -\pi$ and $\theta = +\pi$ map to
the \emph{same point} $(-1, 0, \omega_j/\omega_{\max})$, eliminating the
boundary discontinuity.

\subsection{Spatial RBF Kernel}

We compute pairwise squared Euclidean distances in the cylindrical space:
\begin{equation}
D^2_{(i,j),(k,l)} = \left\|\mathbf{x}_{i,j} - \mathbf{x}_{k,l}\right\|^2
\end{equation}
and form a radial basis function (Gaussian) kernel:
\begin{equation}
K^{\text{spatial}}_{(i,j),(k,l)} = \exp\!\left(-\frac{D^2_{(i,j),(k,l)}}{2\sigma^2}\right)
\end{equation}
with bandwidth $\sigma = 1.0$.  This kernel is symmetric, non-negative,
and equals~1 on the diagonal.

For non-periodic state spaces (e.g.\ Mountain Car), the embedding reduces
to normalised Cartesian coordinates and the kernel becomes a standard
spatial RBF.

\subsection{SR-Based Affinity}

The standard spectral clustering affinity is the symmetrised successor
matrix:
\begin{equation}
\mathbf{A}^{\text{SR}} = \max\!\left(\mathbf{M},\;\mathbf{M}^\top\right)
\end{equation}
To blend it with the spatial kernel, we first normalise
$\mathbf{A}^{\text{SR}}$ to $[0, 1]$:
\begin{equation}
\widetilde{\mathbf{A}}^{\text{SR}} = \frac{\mathbf{A}^{\text{SR}} - \min\!\left(\mathbf{A}^{\text{SR}}\right)}{\max\!\left(\mathbf{A}^{\text{SR}}\right) - \min\!\left(\mathbf{A}^{\text{SR}}\right)}
\end{equation}

\subsection{Adaptive Blending}

A uniform blend weight $\alpha$ applies the spatial kernel equally to all
state pairs, regardless of how well $\mathbf{M}$ has been estimated.
This can be problematic: in environments with walls or barriers, two
states may be spatially adjacent but dynamically disconnected.  A uniform
spatial kernel would incorrectly boost their affinity, blurring the very
bottleneck structure that clustering should discover.

We address this with an \emph{adaptive} blend weight that varies per
state pair based on the confidence in $\mathbf{M}$.  The intuition is
simple: trust $\mathbf{M}$ where it is reliable; fall back on the spatial
kernel where it is not.

\paragraph{Per-state confidence.}
The row sum $r_i = \sum_j M_{ij}$ reflects how well state~$i$ has been
explored: well-visited states accumulate large successor mass, while
rarely-visited states (e.g.\ near-goal boundary regions) have small row
sums.  We normalise to $[0,1]$:
\begin{equation}
\hat{c}_i = \frac{r_i}{\max_k\, r_k}
\end{equation}

\paragraph{Per-pair confidence.}
For the pair $(i,j)$, we take the minimum of the two states'
confidences, so that the spatial kernel activates whenever \emph{either}
state is poorly estimated:
\begin{equation}
c_{ij} = \min\!\left(\hat{c}_i,\;\hat{c}_j\right)
\end{equation}

\paragraph{Adaptive blend weight.}
The effective spatial weight for pair $(i,j)$ is:
\begin{equation}
\alpha_{ij} = \alpha_{\max}\,(1 - c_{ij})
\end{equation}
where $\alpha_{\max}$ is the maximum blend weight (a hyperparameter;
default $0.5$ for the Pendulum, $0.3$ for Mountain Car).  The final
blended affinity becomes:
\begin{equation}
\boxed{A_{ij} = (1 - \alpha_{ij})\,\widetilde{A}^{\text{SR}}_{ij} + \alpha_{ij}\,K^{\text{spatial}}_{ij}}
\end{equation}

This gives the desired behaviour in three regimes:

\begin{itemize}[leftmargin=2em]
\item \textbf{Both states well-visited}
($\hat{c}_i, \hat{c}_j \approx 1$): $\alpha_{ij} \approx 0$.
The affinity is determined almost entirely by the SR---dynamical barriers
(walls, unreachable regions) are preserved.

\item \textbf{At least one state poorly-visited}
($\min(\hat{c}_i, \hat{c}_j) \approx 0$):
$\alpha_{ij} \approx \alpha_{\max}$.  The spatial kernel fills in for the
unreliable SR, preventing sparse boundary states from splitting off as
outlier clusters.

\item \textbf{Periodic boundaries (Pendulum)}:  States at
$\theta \approx \pm\pi$ are typically well-visited (the pendulum swings
through the bottom frequently), so $c_{ij}$ is moderate-to-high.  The
cylindrical kernel still provides a boost because the \emph{spatial}
distance across the branch cut is zero, while the SR may under-represent
this connection due to bin discretisation.  The adaptive weight ensures
this boost is proportionate.
\end{itemize}

The key property is that $\mathbf{K}^{\text{spatial}}$ is invariant to
the choice of angle branch cut---$\theta = -\pi$ and $\theta = +\pi$
have zero distance---so boundary states that are dynamically connected in
$\mathbf{M}$ are no longer artificially separated by the discretisation.

\subsection{Scope of Application}

The blended affinity $\mathbf{A}$ is used exclusively for \emph{macro-state
discovery}: it determines the spectral embedding (for visualisation) and
the spectral clustering (for cluster assignments).  The core micro-level
quantities---the learned successor matrix $\mathbf{M}$, the transition
matrix $\mathbf{B}$, and the value function
$\mathbf{V} = \mathbf{M}\,\mathbf{C}$---remain unmodified.

Macro-level planning matrices ($\mathbf{B}_{\text{macro}}$,
$\mathbf{M}_{\text{macro}}$) are derived from raw transition counts
between the resulting clusters, so they are \emph{indirectly} shaped by
the blended clustering but do not incorporate the spatial kernel
themselves.

This separation is deliberate: the spatial kernel should produce
physically coherent macro-state regions without distorting the learned
dynamics that drive action selection.

\subsection{Related Work}

Machado et~al.~\autocite{machado2018eigenoption} proved that the
eigenvectors of the successor representation matrix $\mathbf{M}$ are
equivalent to proto-value functions---the eigenvectors of the graph
Laplacian over the state-transition graph.  Their \emph{eigenoption
discovery} framework uses spectral decomposition of $\mathbf{M}$ to
identify bottleneck states and construct temporally abstract actions
(options).  The earlier Laplacian framework~\autocite{machado2017laplacian}
showed that eigenoptions correspond to eigenvectors associated with the
largest eigenvalues of the graph Laplacian, providing a principled
connection between spectral methods and option discovery in reinforcement
learning.  Our micro-level spectral clustering on $\mathbf{M}$ follows
this established approach directly.

Clustering data embedded in spaces with periodic boundary conditions is a
well-studied problem in computational science.  The established solution
is to embed periodic coordinates on the unit circle
$(\cos\theta, \sin\theta)$ so that the metric respects the underlying
topology.  This approach has been applied to
$k$-means~\autocite{gao2022kmeans} and
DBSCAN~\autocite{lazzari2025dbscan} in molecular simulation and other
periodic settings.

The idea of combining multiple affinity matrices appears extensively in
multi-view clustering, where a unified affinity is learned from several
``views'' of the same data.  Constrained spectral clustering modifies the
affinity matrix with must-link / cannot-link constraints or spatial
regularisation terms~\autocite{vonluxburg2007tutorial}.  Our convex
combination of a dynamics-based affinity
($\widetilde{\mathbf{A}}^{\text{SR}}$) with a geometry-based kernel
($\mathbf{K}^{\text{spatial}}$) is conceptually analogous, treating the
successor representation and the physical state layout as two
complementary views of the same state space.

To our knowledge, the specific combination of a learned successor
representation with a cylindrical spatial RBF kernel via an adaptive,
confidence-weighted convex blend for macro-state discovery has not
appeared in prior work.  The adaptive confidence weighting is a new
contribution that makes the blending robust to environments containing
dynamical barriers, where a uniform spatial kernel would incorrectly
smooth over true bottleneck structure.

%
% Requires: amsmath, amssymb, enumitem (in parent preamble)
% =====================================================================

\section{Adaptive Spatial--Dynamical Affinity Blending for Macro-State Discovery}

\subsection{Problem}

The pendulum state space is parameterised by $(\theta, \omega)$ where
$\theta \in [-\pi, \pi]$ is \emph{periodic}---the states $\theta = -\pi$
and $\theta = +\pi$ are physically identical (pendulum hanging straight
down).  However, in the discretised grid they occupy opposite edges, so
standard spectral clustering on the successor representation matrix
$\mathbf{M}$ can fragment dynamically-connected boundary states into
separate clusters.

\subsection{Cylindrical Embedding}

Instead of operating in the flat $(\theta, \omega)$ space, we embed each
discrete state $s_{i,j}$ (with bin centres $\theta_i$, $\omega_j$) into
a 3-dimensional \emph{cylindrical coordinate space}:
\begin{equation}
\mathbf{x}_{i,j} = \left(\cos\theta_i,\;\sin\theta_i,\;\frac{\omega_j}{\omega_{\max}}\right)
\end{equation}
where $\omega_{\max} = 8.0\;\text{rad/s}$ is the clipping bound for
angular velocity.  The angle maps onto the unit circle
($|\cos\theta|, |\sin\theta| \leq 1$), and dividing $\omega$ by
$\omega_{\max}$ normalises it to $[-1, 1]$, making all three dimensions
commensurate.

In this embedding, states at $\theta = -\pi$ and $\theta = +\pi$ map to
the \emph{same point} $(-1, 0, \omega_j/\omega_{\max})$, eliminating the
boundary discontinuity.

\subsection{Spatial RBF Kernel}

We compute pairwise squared Euclidean distances in the cylindrical space:
\begin{equation}
D^2_{(i,j),(k,l)} = \left\|\mathbf{x}_{i,j} - \mathbf{x}_{k,l}\right\|^2
\end{equation}
and form a radial basis function (Gaussian) kernel:
\begin{equation}
K^{\text{spatial}}_{(i,j),(k,l)} = \exp\!\left(-\frac{D^2_{(i,j),(k,l)}}{2\sigma^2}\right)
\end{equation}
with bandwidth $\sigma = 1.0$.  This kernel is symmetric, non-negative,
and equals~1 on the diagonal.

For non-periodic state spaces (e.g.\ Mountain Car), the embedding reduces
to normalised Cartesian coordinates and the kernel becomes a standard
spatial RBF.

\subsection{SR-Based Affinity}

The standard spectral clustering affinity is the symmetrised successor
matrix:
\begin{equation}
\mathbf{A}^{\text{SR}} = \max\!\left(\mathbf{M},\;\mathbf{M}^\top\right)
\end{equation}
To blend it with the spatial kernel, we first normalise
$\mathbf{A}^{\text{SR}}$ to $[0, 1]$:
\begin{equation}
\widetilde{\mathbf{A}}^{\text{SR}} = \frac{\mathbf{A}^{\text{SR}} - \min\!\left(\mathbf{A}^{\text{SR}}\right)}{\max\!\left(\mathbf{A}^{\text{SR}}\right) - \min\!\left(\mathbf{A}^{\text{SR}}\right)}
\end{equation}

\subsection{Adaptive Blending}

A uniform blend weight $\alpha$ applies the spatial kernel equally to all
state pairs, regardless of how well $\mathbf{M}$ has been estimated.
This can be problematic: in environments with walls or barriers, two
states may be spatially adjacent but dynamically disconnected.  A uniform
spatial kernel would incorrectly boost their affinity, blurring the very
bottleneck structure that clustering should discover.

We address this with an \emph{adaptive} blend weight that varies per
state pair based on the confidence in $\mathbf{M}$.  The intuition is
simple: trust $\mathbf{M}$ where it is reliable; fall back on the spatial
kernel where it is not.

\paragraph{Per-state confidence.}
The row sum $r_i = \sum_j M_{ij}$ reflects how well state~$i$ has been
explored: well-visited states accumulate large successor mass, while
rarely-visited states (e.g.\ near-goal boundary regions) have small row
sums.  We normalise to $[0,1]$:
\begin{equation}
\hat{c}_i = \frac{r_i}{\max_k\, r_k}
\end{equation}

\paragraph{Per-pair confidence.}
For the pair $(i,j)$, we take the minimum of the two states'
confidences, so that the spatial kernel activates whenever \emph{either}
state is poorly estimated:
\begin{equation}
c_{ij} = \min\!\left(\hat{c}_i,\;\hat{c}_j\right)
\end{equation}

\paragraph{Adaptive blend weight.}
The effective spatial weight for pair $(i,j)$ is:
\begin{equation}
\alpha_{ij} = \alpha_{\max}\,(1 - c_{ij})
\end{equation}
where $\alpha_{\max}$ is the maximum blend weight (a hyperparameter;
default $0.5$ for the Pendulum, $0.3$ for Mountain Car).  The final
blended affinity becomes:
\begin{equation}
\boxed{A_{ij} = (1 - \alpha_{ij})\,\widetilde{A}^{\text{SR}}_{ij} + \alpha_{ij}\,K^{\text{spatial}}_{ij}}
\end{equation}

This gives the desired behaviour in three regimes:

\begin{itemize}[leftmargin=2em]
\item \textbf{Both states well-visited}
($\hat{c}_i, \hat{c}_j \approx 1$): $\alpha_{ij} \approx 0$.
The affinity is determined almost entirely by the SR---dynamical barriers
(walls, unreachable regions) are preserved.

\item \textbf{At least one state poorly-visited}
($\min(\hat{c}_i, \hat{c}_j) \approx 0$):
$\alpha_{ij} \approx \alpha_{\max}$.  The spatial kernel fills in for the
unreliable SR, preventing sparse boundary states from splitting off as
outlier clusters.

\item \textbf{Periodic boundaries (Pendulum)}:  States at
$\theta \approx \pm\pi$ are typically well-visited (the pendulum swings
through the bottom frequently), so $c_{ij}$ is moderate-to-high.  The
cylindrical kernel still provides a boost because the \emph{spatial}
distance across the branch cut is zero, while the SR may under-represent
this connection due to bin discretisation.  The adaptive weight ensures
this boost is proportionate.
\end{itemize}

The key property is that $\mathbf{K}^{\text{spatial}}$ is invariant to
the choice of angle branch cut---$\theta = -\pi$ and $\theta = +\pi$
have zero distance---so boundary states that are dynamically connected in
$\mathbf{M}$ are no longer artificially separated by the discretisation.

\subsection{Scope of Application}

The blended affinity $\mathbf{A}$ is used exclusively for \emph{macro-state
discovery}: it determines the spectral embedding (for visualisation) and
the spectral clustering (for cluster assignments).  The core micro-level
quantities---the learned successor matrix $\mathbf{M}$, the transition
matrix $\mathbf{B}$, and the value function
$\mathbf{V} = \mathbf{M}\,\mathbf{C}$---remain unmodified.

Macro-level planning matrices ($\mathbf{B}_{\text{macro}}$,
$\mathbf{M}_{\text{macro}}$) are derived from raw transition counts
between the resulting clusters, so they are \emph{indirectly} shaped by
the blended clustering but do not incorporate the spatial kernel
themselves.

This separation is deliberate: the spatial kernel should produce
physically coherent macro-state regions without distorting the learned
dynamics that drive action selection.

\subsection{Related Work}

Machado et~al.~\autocite{machado2018eigenoption} proved that the
eigenvectors of the successor representation matrix $\mathbf{M}$ are
equivalent to proto-value functions---the eigenvectors of the graph
Laplacian over the state-transition graph.  Their \emph{eigenoption
discovery} framework uses spectral decomposition of $\mathbf{M}$ to
identify bottleneck states and construct temporally abstract actions
(options).  The earlier Laplacian framework~\autocite{machado2017laplacian}
showed that eigenoptions correspond to eigenvectors associated with the
largest eigenvalues of the graph Laplacian, providing a principled
connection between spectral methods and option discovery in reinforcement
learning.  Our micro-level spectral clustering on $\mathbf{M}$ follows
this established approach directly.

Clustering data embedded in spaces with periodic boundary conditions is a
well-studied problem in computational science.  The established solution
is to embed periodic coordinates on the unit circle
$(\cos\theta, \sin\theta)$ so that the metric respects the underlying
topology.  This approach has been applied to
$k$-means~\autocite{gao2022kmeans} and
DBSCAN~\autocite{lazzari2025dbscan} in molecular simulation and other
periodic settings.

The idea of combining multiple affinity matrices appears extensively in
multi-view clustering, where a unified affinity is learned from several
``views'' of the same data.  Constrained spectral clustering modifies the
affinity matrix with must-link / cannot-link constraints or spatial
regularisation terms~\autocite{vonluxburg2007tutorial}.  Our convex
combination of a dynamics-based affinity
($\widetilde{\mathbf{A}}^{\text{SR}}$) with a geometry-based kernel
($\mathbf{K}^{\text{spatial}}$) is conceptually analogous, treating the
successor representation and the physical state layout as two
complementary views of the same state space.

To our knowledge, the specific combination of a learned successor
representation with a cylindrical spatial RBF kernel via an adaptive,
confidence-weighted convex blend for macro-state discovery has not
appeared in prior work.  The adaptive confidence weighting is a new
contribution that makes the blending robust to environments containing
dynamical barriers, where a uniform spatial kernel would incorrectly
smooth over true bottleneck structure.

%
% Requires: amsmath, amssymb, enumitem (in parent preamble)
% =====================================================================

\section{Adaptive Spatial--Dynamical Affinity Blending for Macro-State Discovery}

\subsection{Problem}

The pendulum state space is parameterised by $(\theta, \omega)$ where
$\theta \in [-\pi, \pi]$ is \emph{periodic}---the states $\theta = -\pi$
and $\theta = +\pi$ are physically identical (pendulum hanging straight
down).  However, in the discretised grid they occupy opposite edges, so
standard spectral clustering on the successor representation matrix
$\mathbf{M}$ can fragment dynamically-connected boundary states into
separate clusters.

\subsection{Cylindrical Embedding}

Instead of operating in the flat $(\theta, \omega)$ space, we embed each
discrete state $s_{i,j}$ (with bin centres $\theta_i$, $\omega_j$) into
a 3-dimensional \emph{cylindrical coordinate space}:
\begin{equation}
\mathbf{x}_{i,j} = \left(\cos\theta_i,\;\sin\theta_i,\;\frac{\omega_j}{\omega_{\max}}\right)
\end{equation}
where $\omega_{\max} = 8.0\;\text{rad/s}$ is the clipping bound for
angular velocity.  The angle maps onto the unit circle
($|\cos\theta|, |\sin\theta| \leq 1$), and dividing $\omega$ by
$\omega_{\max}$ normalises it to $[-1, 1]$, making all three dimensions
commensurate.

In this embedding, states at $\theta = -\pi$ and $\theta = +\pi$ map to
the \emph{same point} $(-1, 0, \omega_j/\omega_{\max})$, eliminating the
boundary discontinuity.

\subsection{Spatial RBF Kernel}

We compute pairwise squared Euclidean distances in the cylindrical space:
\begin{equation}
D^2_{(i,j),(k,l)} = \left\|\mathbf{x}_{i,j} - \mathbf{x}_{k,l}\right\|^2
\end{equation}
and form a radial basis function (Gaussian) kernel:
\begin{equation}
K^{\text{spatial}}_{(i,j),(k,l)} = \exp\!\left(-\frac{D^2_{(i,j),(k,l)}}{2\sigma^2}\right)
\end{equation}
with bandwidth $\sigma = 1.0$.  This kernel is symmetric, non-negative,
and equals~1 on the diagonal.

For non-periodic state spaces (e.g.\ Mountain Car), the embedding reduces
to normalised Cartesian coordinates and the kernel becomes a standard
spatial RBF.

\subsection{SR-Based Affinity}

The standard spectral clustering affinity is the symmetrised successor
matrix:
\begin{equation}
\mathbf{A}^{\text{SR}} = \max\!\left(\mathbf{M},\;\mathbf{M}^\top\right)
\end{equation}
To blend it with the spatial kernel, we first normalise
$\mathbf{A}^{\text{SR}}$ to $[0, 1]$:
\begin{equation}
\widetilde{\mathbf{A}}^{\text{SR}} = \frac{\mathbf{A}^{\text{SR}} - \min\!\left(\mathbf{A}^{\text{SR}}\right)}{\max\!\left(\mathbf{A}^{\text{SR}}\right) - \min\!\left(\mathbf{A}^{\text{SR}}\right)}
\end{equation}

\subsection{Adaptive Blending}

A uniform blend weight $\alpha$ applies the spatial kernel equally to all
state pairs, regardless of how well $\mathbf{M}$ has been estimated.
This can be problematic: in environments with walls or barriers, two
states may be spatially adjacent but dynamically disconnected.  A uniform
spatial kernel would incorrectly boost their affinity, blurring the very
bottleneck structure that clustering should discover.

We address this with an \emph{adaptive} blend weight that varies per
state pair based on the confidence in $\mathbf{M}$.  The intuition is
simple: trust $\mathbf{M}$ where it is reliable; fall back on the spatial
kernel where it is not.

\paragraph{Per-state confidence.}
The row sum $r_i = \sum_j M_{ij}$ reflects how well state~$i$ has been
explored: well-visited states accumulate large successor mass, while
rarely-visited states (e.g.\ near-goal boundary regions) have small row
sums.  We normalise to $[0,1]$:
\begin{equation}
\hat{c}_i = \frac{r_i}{\max_k\, r_k}
\end{equation}

\paragraph{Per-pair confidence.}
For the pair $(i,j)$, we take the minimum of the two states'
confidences, so that the spatial kernel activates whenever \emph{either}
state is poorly estimated:
\begin{equation}
c_{ij} = \min\!\left(\hat{c}_i,\;\hat{c}_j\right)
\end{equation}

\paragraph{Adaptive blend weight.}
The effective spatial weight for pair $(i,j)$ is:
\begin{equation}
\alpha_{ij} = \alpha_{\max}\,(1 - c_{ij})
\end{equation}
where $\alpha_{\max}$ is the maximum blend weight (a hyperparameter;
default $0.5$ for the Pendulum, $0.3$ for Mountain Car).  The final
blended affinity becomes:
\begin{equation}
\boxed{A_{ij} = (1 - \alpha_{ij})\,\widetilde{A}^{\text{SR}}_{ij} + \alpha_{ij}\,K^{\text{spatial}}_{ij}}
\end{equation}

This gives the desired behaviour in three regimes:

\begin{itemize}[leftmargin=2em]
\item \textbf{Both states well-visited}
($\hat{c}_i, \hat{c}_j \approx 1$): $\alpha_{ij} \approx 0$.
The affinity is determined almost entirely by the SR---dynamical barriers
(walls, unreachable regions) are preserved.

\item \textbf{At least one state poorly-visited}
($\min(\hat{c}_i, \hat{c}_j) \approx 0$):
$\alpha_{ij} \approx \alpha_{\max}$.  The spatial kernel fills in for the
unreliable SR, preventing sparse boundary states from splitting off as
outlier clusters.

\item \textbf{Periodic boundaries (Pendulum)}:  States at
$\theta \approx \pm\pi$ are typically well-visited (the pendulum swings
through the bottom frequently), so $c_{ij}$ is moderate-to-high.  The
cylindrical kernel still provides a boost because the \emph{spatial}
distance across the branch cut is zero, while the SR may under-represent
this connection due to bin discretisation.  The adaptive weight ensures
this boost is proportionate.
\end{itemize}

The key property is that $\mathbf{K}^{\text{spatial}}$ is invariant to
the choice of angle branch cut---$\theta = -\pi$ and $\theta = +\pi$
have zero distance---so boundary states that are dynamically connected in
$\mathbf{M}$ are no longer artificially separated by the discretisation.

\subsection{Scope of Application}

The blended affinity $\mathbf{A}$ is used exclusively for \emph{macro-state
discovery}: it determines the spectral embedding (for visualisation) and
the spectral clustering (for cluster assignments).  The core micro-level
quantities---the learned successor matrix $\mathbf{M}$, the transition
matrix $\mathbf{B}$, and the value function
$\mathbf{V} = \mathbf{M}\,\mathbf{C}$---remain unmodified.

Macro-level planning matrices ($\mathbf{B}_{\text{macro}}$,
$\mathbf{M}_{\text{macro}}$) are derived from raw transition counts
between the resulting clusters, so they are \emph{indirectly} shaped by
the blended clustering but do not incorporate the spatial kernel
themselves.

This separation is deliberate: the spatial kernel should produce
physically coherent macro-state regions without distorting the learned
dynamics that drive action selection.

\subsection{Related Work}

Machado et~al.~\autocite{machado2018eigenoption} proved that the
eigenvectors of the successor representation matrix $\mathbf{M}$ are
equivalent to proto-value functions---the eigenvectors of the graph
Laplacian over the state-transition graph.  Their \emph{eigenoption
discovery} framework uses spectral decomposition of $\mathbf{M}$ to
identify bottleneck states and construct temporally abstract actions
(options).  The earlier Laplacian framework~\autocite{machado2017laplacian}
showed that eigenoptions correspond to eigenvectors associated with the
largest eigenvalues of the graph Laplacian, providing a principled
connection between spectral methods and option discovery in reinforcement
learning.  Our micro-level spectral clustering on $\mathbf{M}$ follows
this established approach directly.

Clustering data embedded in spaces with periodic boundary conditions is a
well-studied problem in computational science.  The established solution
is to embed periodic coordinates on the unit circle
$(\cos\theta, \sin\theta)$ so that the metric respects the underlying
topology.  This approach has been applied to
$k$-means~\autocite{gao2022kmeans} and
DBSCAN~\autocite{lazzari2025dbscan} in molecular simulation and other
periodic settings.

The idea of combining multiple affinity matrices appears extensively in
multi-view clustering, where a unified affinity is learned from several
``views'' of the same data.  Constrained spectral clustering modifies the
affinity matrix with must-link / cannot-link constraints or spatial
regularisation terms~\autocite{vonluxburg2007tutorial}.  Our convex
combination of a dynamics-based affinity
($\widetilde{\mathbf{A}}^{\text{SR}}$) with a geometry-based kernel
($\mathbf{K}^{\text{spatial}}$) is conceptually analogous, treating the
successor representation and the physical state layout as two
complementary views of the same state space.

To our knowledge, the specific combination of a learned successor
representation with a cylindrical spatial RBF kernel via an adaptive,
confidence-weighted convex blend for macro-state discovery has not
appeared in prior work.  The adaptive confidence weighting is a new
contribution that makes the blending robust to environments containing
dynamical barriers, where a uniform spatial kernel would incorrectly
smooth over true bottleneck structure.
